\documentclass[../../../include/open-logic-section]{subfiles}

\begin{document}

\olfileid{sth}{replacement}{limofsize}
\olsection{التقييد بالحجم}

أشهر ما يسلك في تسويغ الاستبدال تسويغًا «داخليًا» أن يستند إلى
المبدأ الآتي:
\begin{enumerate}
	\item[] \limofsize. تؤلف أي أشياء مجموعة، شريطة ألا يكون عددها
	كبيرًا أكثر مما ينبغي.
\end{enumerate}
وهذا المبدأ يفضي إلى الاستبدال رأسًا. فكل مجموعة يحدثها الاستبدال صورة
شاملة للمجموعة التي صدرت عنها. ولنضبط القول: إذا أنشأنا بالاستبدال
$\funimage{\tau}{\OLId{latin-looped}{A}} = \Setabs{\tau(\OLId{latin-ordinary}{x})}{\OLId{latin-ordinary}{x} \in \OLId{latin-looped}{A}}$، كانت القاعدة
$\tau$ مرسلة كل~!!{element} من~$\OLId{latin-looped}{A}$ إلى صورته، وكانت شاملة على
$\funimage{\tau}{\OLId{latin-looped}{A}}$. فعدد~$\funimage{\tau}{\OLId{latin-looped}{A}}$ لا يزيد، بمعيار
الدوال الشاملة، على عدد~$\OLId{latin-looped}{A}$. ولا نستخرج من ذلك، في العموم، المتباينة
الكاردينالية الأقوى ما لم نضف مسلّمة الاختيار. فإذا كانت عناصر~$\OLId{latin-looped}{A}$
غير كثيرة إلى حد يمنعها من تأليف مجموعة، لم تكن صورها كثيرة إلى حد
يمنع تأليف~$\funimage{\tau}{\OLId{latin-looped}{A}}$.

غير أن أول عائق في اتخاذ~\limofsize{} حجة للاستبدال هو أننا لم نضع
إلى الآن مبدأ من هذا القبيل. بل إن روايتنا للتصور التراكمي-التكراري
في~\crefrange{sth:story::chap}{sth:z::chap} لم تومئ إليه قط. ولهذا
قال بولوس في موضع: «لعل المرء يستطيع أن يستنتج أن وراء نظرية
المجموعات فكرتين على الأقل»~\citeyearpar[p.~19]{Boolos1989}. فالتصور
التراكمي-التكراري يراد به تسويغ~$\Z$، وأما~\limofsize{} فيراد به
تسويغ الاستبدال.

وليس الأمر مجرد سكوتنا السابق عن~\limofsize؛ بل إن المبدأ بصورته
هذه يبدو بعيدًا عن الانسجام مع التصور التراكمي-التكراري، إذ لا يذكر
شيئًا من تكون المجموعات في \emph{مراحل}.

ولا غرابة في ذلك إذا نظرنا إلى تاريخ هاتين الفكرتين، فإنهما ترجعان
إلى مشروعين مختلفين لدفع مفارقات نظرية المجموعات. فالتصور
التراكمي-التكراري يدفع مفارقة راسل، في الجملة، بأن مجموعة راسل لم يكن
منتظرًا وجودها لأنها لا تتكون في مرحلة قط. وأما~\limofsize{} فيدفعها
بأن تلك المجموعة لم يكن منتظرًا وجودها لأنها أكبر من أن توجد.

ومتى صيغ الفرق على هذا الوجه ظهر أن~\limofsize{}، بوصفه جوابًا عن
المفارقات، يحتاج إلى حجة أقوى بكثير. تأمل مثلًا هذه الصورة من مفارقة
راسل: \emph{لا كلب قصير الأنف يشم بالضبط جميع الكلاب قصيرة الأنف
التي لا تشم نفسها} (انظر~\olref[sth][story][rus]{sec}). لو سأل سائل:
«لم لا يوجد هذا الكلب؟» لم يكن جوابًا مقبولًا أن يقال: لأنه سيضطر إلى
شم عدد مفرط من الكلاب. فكيف يصير القول إن مجموعة راسل «أكبر من أن
توجد» تفسيرًا حدسيًا حسنًا لعدم وجودها؟

والحاصل أن بقاء المرء متحيرًا في الدافع «الحدسي» لـ\limofsize{} أمر
له عذره.

\end{document}
